\phantomsection
\section*{结 \quad 论}
\addcontentsline{toc}{section}{结论}


本文面向 MNIST 手写数字识别任务，设计了一种参数量约 0.47M 的小型卷积神经网络，并系统研究了数据增强、批量归一化、Dropout 与自适应优化等训练方法。实验表明，所提方法在测试集上取得 99.32\% 的识别准确率，优于 Softmax 回归、多层感知机与 LeNet-5 等对比方法，消融实验验证了各训练技巧的独立贡献。


本文方法仍存在不足：实验仅在 MNIST 这一相对规整的数据集上验证，对更复杂的真实手写场景，其泛化能力有待进一步检验；网络结构与训练策略也有优化空间。未来可在更具挑战性的数据集上开展实验，并探索残差连接、注意力机制以及模型压缩等方向。
